\documentclass[11pt,a4paper]{article}
\usepackage[utf8]{inputenc}
\usepackage[T1]{fontenc}
\usepackage{amsmath, amssymb, amsthm, mathrsfs, mathtools}
\usepackage[margin=1in]{geometry}
\usepackage{enumerate}
\usepackage{hyperref}
\usepackage{xcolor}
\hypersetup{
  colorlinks=true,
  linkcolor=blue!50!black,
  citecolor=blue!50!black,
  urlcolor=blue!50!black
}

%%% Theorem environments
\theoremstyle{plain}
\newtheorem{theorem}{Theorem}[section]
\newtheorem{proposition}[theorem]{Proposition}
\newtheorem{lemma}[theorem]{Lemma}
\newtheorem{corollary}[theorem]{Corollary}

\theoremstyle{definition}
\newtheorem{definition}[theorem]{Definition}
\newtheorem{example}[theorem]{Example}

\theoremstyle{remark}
\newtheorem{remark}[theorem]{Remark}

%%% Macros
\newcommand{\C}{\mathbb{C}}
\newcommand{\R}{\mathbb{R}}
\newcommand{\Z}{\mathbb{Z}}
\newcommand{\Q}{\mathbb{Q}}
\newcommand{\PGL}{\operatorname{PGL}}
\newcommand{\Aut}{\operatorname{Aut}}
\newcommand{\Gal}{\operatorname{Gal}}
\newcommand{\rk}{\operatorname{rank}}
\newcommand{\F}{\mathcal{F}}
\renewcommand{\d}{\,\mathrm{d}}

\title{\textbf{From Goursat's Pseudo-Elliptic Integrals\\
to a Cube-Root Generalization}}
\author{}
\date{\today}

\begin{document}
\maketitle

\begin{abstract}
We give a self-contained, modern exposition of \'Emile Goursat's 1887 theorem on pseudo-elliptic integrals --- those integrals of the form $\int F(t)\,\d t/\sqrt{R(t)}$ with $R$ a cubic or quartic polynomial that, despite living on a genus-$1$ algebraic curve, admit elementary antiderivatives. After reviewing integration in finite terms and Liouville's theorem, we present Goursat's two main theorems with proofs phrased in the language of M\"obius automorphisms of the underlying hyperelliptic curve. We then develop a cube-root analog: for integrals of the form $\int F(t)\,\d t/\sqrt[3]{R(t)}$ with $R$ cubic, an order-$3$ M\"obius substitution cyclically permuting the roots of $R$ induces an eigendecomposition into three pieces. Two of the three eigenpieces (eigenvalues $1$ and $\omega^2$, where $\omega = e^{2\pi i/3}$) descend through a chain of substitutions to genus-$0$ curves and yield elementary antiderivatives; the middle eigenpiece (eigenvalue $\omega$) descends only to the genus-$1$ curve $y^3 = x(x-K)$ and is generically transcendental. We prove this dichotomy and illustrate it with several concrete computations, including the elementary $\int \d t/\sqrt[3]{t^3-1}$ and the genuinely transcendental $\int t\,\d t/\sqrt[3]{t^3-1}$.
\end{abstract}

\tableofcontents

\section{Introduction}
\label{sec:intro}

Among the very first surprises of integral calculus is the asymmetry between
\[
  \int \frac{\d t}{\sqrt{1-t^2}} = \arcsin t + C
\qquad\text{and}\qquad
  \int \frac{\d t}{\sqrt{(1-t^2)(1-k^2t^2)}} = ???
\]
The first integral is elementary; the second is not. Both integrands look formally similar --- a rational function divided by the square root of a polynomial --- yet only the first admits a closed-form antiderivative in terms of the standard library of elementary functions. The deep reason is geometric: the curve $y^2 = 1 - t^2$ has \emph{genus zero} (it is the unit circle), while the curve $y^2 = (1-t^2)(1-k^2 t^2)$ has \emph{genus one}, and on a curve of positive genus most $1$-forms have no elementary antiderivative. The classical incomplete elliptic integral on the right is the prototype.

This raises a refined question: can a $1$-form on a positive-genus curve \emph{ever} integrate elementarily? The answer is yes, but only for very special integrands. The \emph{pseudo-elliptic integrals} are precisely those of the form $\int F(t)\,\d t/\sqrt{R(t)}$ which look elliptic but evaluate elementarily. Their existence is non-obvious, and producing systematic methods for detecting and reducing them was a recurrent theme in nineteenth-century analysis. \'Emile Goursat's note of 1887~\cite{Goursat1887} is the most elegant general account.

The purpose of this paper is twofold. First, in Sections~\ref{sec:liouville}--\ref{sec:goursat}, we give a self-contained modern exposition of Goursat's two main theorems. The framework is standard differential algebra: an integral is elementary precisely when it satisfies the structural conclusion of Liouville's theorem~\cite{Rosenlicht1972}, and Goursat's hypothesis turns out to be a clean Galois-theoretic condition imposed by the M\"obius automorphisms of the curve $y^2 = R(t)$. Second, in Sections~\ref{sec:cuberoot}--\ref{sec:examples}, we develop a cube-root analog of Goursat's theory. The natural setting is the curve $y^3 = R(t)$, and we prove an analog of Goursat's first theorem in which the binary anti-invariance condition $F(t)+F(S(t))=0$ is replaced by an eigenvalue condition under an order-$3$ M\"obius substitution.

The cube-root case exhibits a structural feature with no analog in the square-root world: of the three eigenspaces under a cyclic-Möbius action, two lead to elementary antiderivatives but the middle one stalls at a \emph{cube-root elliptic curve} $y^3=x(x-K)$. This curve also has genus one, and its standard holomorphic differential is generically transcendental. The result is a clean sufficient condition for a cube-root integral to be elementary, together with a precise description of the obstruction when the condition fails.

\medskip

The structure of the paper is as follows. Section~\ref{sec:liouville} reviews integration in finite terms and states Liouville's principle in the form we will use. Section~\ref{sec:goursat} is the modern presentation of Goursat: Theorem~\ref{thm:goursat1} (single-involution) and Theorem~\ref{thm:goursat2} (Klein four-group sum). Section~\ref{sec:cuberoot} develops the cube-root analog: setup, statement of Theorem~\ref{thm:cubic-goursat}, proof, and a discussion of the middle-eigenvalue obstruction. Section~\ref{sec:examples} works through six representative examples. Section~\ref{sec:remarks} closes with remarks on extensions and the relation to differential Galois theory.


\section{Integration in finite terms and Liouville's theorem}
\label{sec:liouville}

The framework for making ``elementary integration'' precise is the differential algebra of Joseph Liouville and its modern reformulation by Ritt~\cite{Ritt1948} and Rosenlicht~\cite{Rosenlicht1972}.

\begin{definition}\label{def:diffield}
A \emph{differential field} is a field $K$ of characteristic zero equipped with a derivation, i.e.\ a map $D \colon K \to K$ satisfying $D(a+b)=Da+Db$ and $D(ab)=aDb+bDa$. We write $a' = Da$ throughout.
\end{definition}

The basic example is $K = \C(t)$ with $D = \d/\d t$. We will also consider algebraic extensions $K\subseteq L$, where the derivation extends uniquely.

\begin{definition}\label{def:elementary}
Let $K\subseteq L$ be differential fields. We say $L$ is an \emph{elementary extension} of $K$ if there is a tower
\[
  K = K_0 \subseteq K_1 \subseteq K_2 \subseteq \cdots \subseteq K_n = L
\]
such that each $K_{i+1} = K_i(\theta_i)$ where one of the following holds:
\begin{enumerate}[(i)]
  \item $\theta_i$ is \emph{algebraic} over $K_i$;
  \item $\theta_i$ is an \emph{exponential}: $\theta_i'/\theta_i \in K_i$;
  \item $\theta_i$ is a \emph{logarithm}: $\theta_i' = \eta'/\eta$ for some $\eta \in K_i$.
\end{enumerate}
An element $f \in K$ is said to have an \emph{elementary integral} (or \emph{integrate in finite terms}) if there exists an elementary extension $L \supseteq K$ and an element $g\in L$ with $g' = f$.
\end{definition}

The class of elementary functions in the calculus textbook sense corresponds to elementary extensions of $\C(t)$.

\begin{theorem}[Liouville's Theorem; \cite{Rosenlicht1972}]\label{thm:liouville}
Let $K$ be a differential field with field of constants $\C$. If $f\in K$ has an elementary integral, then there exist $v\in K$, constants $c_1,\dots,c_n\in \C$, and elements $u_1,\dots,u_n \in K$ such that
\begin{equation}\label{eq:liouville}
  f \;=\; v' \;+\; \sum_{i=1}^n c_i\,\frac{u_i'}{u_i}.
\end{equation}
\end{theorem}

The remarkable feature of \eqref{eq:liouville} is that no new transcendentals appear: although the antiderivative may live in an extension $L\supsetneq K$, the structural form of $f$ is dictated by elements of $K$ itself. This is the precise statement that ``the integral, if elementary, is the derivative of a rational function plus a sum of constant multiples of logarithms of rational functions.''

For our purposes we need a refinement which permits the field $K$ to be an algebraic extension of a rational function field.

\begin{theorem}[Liouville's Theorem for algebraic extensions]\label{thm:liouville-alg}
Let $K = \C(t)(\theta)$ where $\theta$ is algebraic over $\C(t)$ of degree $n$, and let $f\in K$. If $\int f\,\d t$ is elementary, then it can be written
\[
  \int f\,\d t = v + \sum_{i=1}^m c_i \log u_i
\]
for some $v, u_1,\dots,u_m \in K$ and constants $c_i \in \C$.
\end{theorem}

This is the form one applies to integrands of the type $F(t)/\sqrt{R(t)}$ or $F(t)/\sqrt[3]{R(t)}$. The relevant field is $K = \C(t,\theta)$ where $\theta^2 = R(t)$ (respectively $\theta^3 = R(t)$).

\begin{remark}
Theorem~\ref{thm:liouville-alg} has powerful consequences. For instance, if $X$ is a smooth projective curve of genus $g\geq 1$ and $\omega$ is a holomorphic $1$-form on $X$ (a so-called \emph{differential of the first kind}), then the function $P\mapsto \int_{P_0}^P \omega$ is \emph{never} elementary as a function on $X$~\cite{Bronstein2005}. This is a deep refinement of the classical fact that $\int \d t/\sqrt{(1-t^2)(1-k^2t^2)}$ is not elementary, and we will rely on it in Section~\ref{sec:cuberoot} to identify the obstruction in the cube-root case.
\end{remark}

\begin{remark}
The constructive content of Theorem~\ref{thm:liouville} is the algorithm of Risch~\cite{Risch1969}, refined for algebraic integrands by Trager~\cite{Trager1984} and exposited in Bronstein~\cite{Bronstein2005}. Goursat's theorems below predate these developments by nearly a century and provide closed-form reductions in cases where the algebraic structure is particularly transparent.
\end{remark}


\section{Goursat's pseudo-elliptic theorem}
\label{sec:goursat}

Throughout this section, $R(t)\in\C[t]$ denotes a polynomial of degree $3$ or $4$ with \emph{simple roots} — i.e., all roots distinct, or equivalently $\gcd(R,R') = 1$. We write $C := \{(t,y)\in\overline{\C^2} : y^2 = R(t)\}$ for the smooth projective completion of the affine curve $y^2=R(t)$; this is a curve of genus~$1$. We are interested in the integrals
\begin{equation}\label{eq:elliptic-integrand}
  I = \int F(t)\,\frac{\d t}{\sqrt{R(t)}}
\end{equation}
with $F\in\C(t)$.

\begin{definition}\label{def:pseudo-elliptic}
The integral~\eqref{eq:elliptic-integrand} is called \emph{pseudo-elliptic} if it is elementary in the sense of Definition~\ref{def:elementary}, i.e.\ if it admits a closed-form antiderivative built from rational functions, $\sqrt{R}$, exponentials, and logarithms.
\end{definition}

The reason for the term \emph{pseudo-elliptic} is that the integrand lives most naturally as a $1$-form on $C$, and a generic such integral is genuinely transcendental (an elliptic integral). The pseudo-elliptic ones are the rare exceptions where the integral happens to be elementary.

\subsection{The role of M\"obius involutions}

The key idea is to exploit symmetries of the curve $C$. The M\"obius group $\PGL_2(\C)$ acts on $\C\cup\{\infty\}=\mathbb{P}^1(\C)$ by fractional linear transformations $S(t) = (At+B)/(Ct+D)$. An element $S \in \PGL_2(\C)$ is an \emph{involution} if $S^2 = \mathrm{id}$ and $S\neq \mathrm{id}$. A non-trivial M\"obius involution has exactly two fixed points on $\mathbb{P}^1(\C)$.

\begin{lemma}\label{lem:involutions-pairs}
A non-trivial M\"obius involution $S$ on $\mathbb{P}^1(\C)$ permuting the roots of a quartic $R$ with simple roots must permute them in two pairs (no $S$-fixed root).

If the four roots are $\{a,b,c,d\}$ and $S$ swaps $\{a,b\}$ with $\{c,d\}$, then explicitly
\begin{equation}\label{eq:involution-formula}
  S(t) \;=\; \frac{(ab-cd)\,t + (a+b)cd - (c+d)ab}{[(a+b)-(c+d)]\,t - (ab-cd)}.
\end{equation}
There are exactly three such involutions, corresponding to the three ways of pairing four points; together with the identity they form a Klein four-group $V_4 \subseteq \PGL_2(\C)$.
\end{lemma}

\begin{proof}
A M\"obius transformation is determined by its action on three points; once the values $S(a)=b$, $S(b)=a$, $S(c)=d$ are fixed (and consistency requires $S(d)=c$), the transformation is uniquely determined and one verifies~\eqref{eq:involution-formula} by direct computation. The product of any two of these three involutions is the third (as one checks on the pairings), so they generate a $V_4$ subgroup.
\end{proof}

\subsection{Goursat's first theorem}

\begin{theorem}[Goursat~\cite{Goursat1887}, Theorem 1]\label{thm:goursat1}
Let $R\in\C[t]$ be a polynomial of degree $3$ or $4$ with simple roots, and let $S\in\PGL_2(\C)$ be a non-trivial M\"obius involution permuting the roots of $R$ in pairs. If $F\in\C(t)$ satisfies the \emph{anti-invariance} condition
\begin{equation}\label{eq:antiinv}
  F(t) + F(S(t)) = 0\qquad \text{(identically)},
\end{equation}
then the integral $\int F(t)\,\d t/\sqrt{R(t)}$ is pseudo-elliptic. Explicitly, let $\alpha,\beta\in\mathbb{P}^1(\C)$ be the two fixed points of $S$, and set
\[
  u = \frac{t-\alpha}{t-\beta},\qquad x = u^2.
\]
Then
\[
  \int F(t)\,\frac{\d t}{\sqrt{R(t)}} \;=\; \int \frac{G(x)\,\d x}{2\sqrt{Q(x)}}
\]
for some $G\in\C(x)$ and quadratic $Q\in\C[x]$, the latter integral being elementary.
\end{theorem}

\begin{proof}
We argue in five steps.

\smallskip\noindent
\textbf{Step 1 (canonical form of $S$).}
Since $S$ has order $2$ with two distinct fixed points $\alpha,\beta$, the M\"obius transformation $u = (t-\alpha)/(t-\beta)$ sends $\alpha\mapsto 0$, $\beta\mapsto\infty$, and conjugates $S$ into a M\"obius transformation of $u$ fixing $\{0,\infty\}$. Such a transformation is multiplication by a scalar $\lambda$, and the involution condition $\lambda^2=1$ together with $\lambda\neq 1$ forces $\lambda = -1$. Thus in $u$-coordinates, $S$ acts as $u\mapsto -u$. We write $t(u) = (\alpha-\beta u)/(1-u)$ for the inverse change of variable.

\smallskip\noindent
\textbf{Step 2 (transformation of $R$).}
Each linear factor $t-r_i$ becomes
\[
  t(u) - r_i \;=\; \frac{(\alpha-r_i) - u(\beta-r_i)}{1-u}.
\]
Hence (assuming $R$ is quartic; the cubic case is identical with $r_4=\infty$ and one factor reading $1$)
\[
  R(t(u))(1-u)^4 \;=\; \prod_{i=1}^4 \big[(\alpha-r_i) - u(\beta-r_i)\big].
\]
The roots $u_i$ of this polynomial in $u$ are $u_i = (\alpha-r_i)/(\beta-r_i) = -(r_i-\alpha)/(\beta-r_i) \cdot \ldots$, more usefully, the image of $r_i$ under $r\mapsto (r-\alpha)/(r-\beta)$. Since $S$ swaps the roots $r_i$ in pairs and acts as $u\mapsto -u$, the $u_i$ come in pairs $\{u_1,-u_1\}$ and $\{u_3,-u_3\}$. Therefore
\begin{equation}\label{eq:R-becomes-even}
  R(t(u))(1-u)^4 \;=\; C\,(u^2 - u_1^2)(u^2 - u_3^2)
\end{equation}
for some non-zero constant $C$. The right-hand side is a polynomial in $u^2$.

\smallskip\noindent
\textbf{Step 3 (transformation of the differential).}
A direct computation gives
\[
  \d t = \frac{\alpha-\beta}{(1-u)^2}\,\d u.
\]
Combined with~\eqref{eq:R-becomes-even} this yields
\begin{equation}\label{eq:differential-becomes}
  \frac{\d t}{\sqrt{R(t)}} \;=\; \frac{(\alpha-\beta)\,\d u/(1-u)^2}{\sqrt{C\,(u^2-u_1^2)(u^2-u_3^2)}/(1-u)^2}
  \;=\; \frac{(\alpha-\beta)\,\d u}{\sqrt{C\,(u^2-u_1^2)(u^2-u_3^2)}}.
\end{equation}
The factors $(1-u)^2$ from the Jacobian and from the radical exactly cancel.

\smallskip\noindent
\textbf{Step 4 (oddness of $F$ in $u$).}
Anti-invariance~\eqref{eq:antiinv} reads, in $u$-coordinates, $F(t(-u)) = F(S(t(u))) = -F(t(u))$. Hence $u\mapsto F(t(u))$ is an odd rational function of $u$, and admits a representation
\[
  F(t(u)) \;=\; u\,G(u^2)
\]
for some $G\in\C(x)$.

\smallskip\noindent
\textbf{Step 5 (descent via $x = u^2$).}
Combining the previous steps:
\[
  \int F(t)\,\frac{\d t}{\sqrt{R(t)}}
  \;=\; (\alpha-\beta)\,\int \frac{u\,G(u^2)\,\d u}{\sqrt{C\,(u^2-u_1^2)(u^2-u_3^2)}}.
\]
Substituting $x=u^2$, $\d x = 2u\,\d u$:
\[
  \;=\; \frac{\alpha-\beta}{2\sqrt{C}}\,\int \frac{G(x)\,\d x}{\sqrt{(x-u_1^2)(x-u_3^2)}}.
\]
The radical is now $\sqrt{Q(x)}$ for the quadratic $Q(x)=(x-u_1^2)(x-u_3^2)$. The integral $\int G(x)\,\d x/\sqrt{Q(x)}$ with $G$ rational and $Q$ quadratic is elementary by classical methods (partial fractions of $G$ together with Euler substitution). This completes the proof.
\end{proof}

\begin{remark}
The case where one of the fixed points is at infinity (so that $S(t) = 2\alpha - t$) is covered by the same argument, the substitution simplifying to $u = t-\alpha$.
\end{remark}

\subsection{Goursat's second theorem}

The three involutions of Lemma~\ref{lem:involutions-pairs} together with the identity form $V_4$, and one can ask for a sum-of-orbits criterion that detects pseudo-ellipticity.

\begin{theorem}[Goursat~\cite{Goursat1887}, Theorem 2]\label{thm:goursat2}
Let $R\in\C[t]$ be a polynomial of degree $3$ or $4$ with simple roots, with M\"obius involutions $S_1,S_2,S_3$ permuting the roots in pairs, and let $F\in\C(t)$. If
\begin{equation}\label{eq:V4-sum}
  F + F\circ S_1 + F\circ S_2 + F\circ S_3 \;=\; 0
\end{equation}
identically, then $\int F\,\d t/\sqrt{R}$ is pseudo-elliptic.
\end{theorem}

\begin{proof}
The group $V_4=\{I,S_1,S_2,S_3\}$ has four irreducible characters $\chi_0,\chi_1,\chi_2,\chi_3 \colon V_4\to\{\pm 1\}$, with $\chi_0$ the trivial character. We label them so that $\chi_j(S_j) = +1$ for $j=1,2,3$, and $\chi_j(S_k) = -1$ for $j\neq k$ in $\{1,2,3\}$.

The associated character projections decompose $\C(t)$ as a direct sum of $\chi_j$-eigenspaces under the natural $V_4$-action. Setting
\[
  F^{(j)} \;:=\; \frac{1}{4}\sum_{g\in V_4} \chi_j(g)\,F\circ g,
\]
hypothesis~\eqref{eq:V4-sum} says exactly $F^{(0)} = 0$, hence $F = F^{(1)} + F^{(2)} + F^{(3)}$.

We claim $F^{(j)}$ is anti-invariant under $S_k$ for every $k\neq j$ (and invariant under $S_j$, as is automatic). Indeed, applying $S_k$ to $F^{(j)}$ permutes the summands by the action $g\mapsto gS_k$, which on characters acts by $\chi_j(g)\mapsto \chi_j(gS_k) = \chi_j(g)\chi_j(S_k)$. For $k\neq j$ this multiplies the projection by $-1$, giving anti-invariance.

By Theorem~\ref{thm:goursat1}, applied to $F^{(j)}$ with any involution $S_k$, $k\neq j$, each integral $\int F^{(j)}\,\d t/\sqrt{R}$ is pseudo-elliptic. Summing the three antiderivatives yields an elementary expression for $\int F\,\d t/\sqrt{R}$.
\end{proof}

\begin{remark}
The character labelling above corrects a common indexing issue. Goursat's original argument writes the projections as
\begin{align*}
  F^{(1)} &= (F - F\!\circ\! S_1 - F\!\circ\! S_2 + F\!\circ\! S_3)/4, \\
  F^{(2)} &= (F - F\!\circ\! S_1 + F\!\circ\! S_2 - F\!\circ\! S_3)/4, \\
  F^{(3)} &= (F + F\!\circ\! S_1 - F\!\circ\! S_2 - F\!\circ\! S_3)/4,
\end{align*}
and one must check carefully which involution $S_k$ to pair with each $F^{(j)}$ when invoking Theorem~\ref{thm:goursat1}: $F^{(j)}$ is anti-invariant under exactly the two involutions $S_k$ with $k\neq j$, hence either of those two works.
\end{remark}

\subsection{Geometric interpretation}

Although Goursat does not phrase his theorems in this way, the modern reading is illuminating. Each M\"obius involution $S$ of $\mathbb{P}^1$ permuting the roots of $R$ extends to an involution $\widetilde{S}$ of the curve $C \colon y^2 = R(t)$ that sends $y$ to $\pm y \cdot J_S(t)$ for an explicit Jacobian factor $J_S$. The pseudo-elliptic condition is then the statement that the differential $F\,\d t/y$ is a $(-1)$-eigenvector of $\widetilde{S}^*$ acting on $\Omega^1_C$. Theorem~\ref{thm:goursat2}'s sum criterion is the orthogonality of $F\,\d t/y$ to the trivial-character isotypic component of $\Omega^1_C$ under the $V_4$-action --- so all of $F\,\d t/y$ lives in non-trivial isotypics, each of which is an eigenspace of \emph{some} involution.

We will see in the next section that the analogous picture for cube radicals is more subtle: there are three eigenspaces under a $\Z/3$-action, but only two of them yield elementary integrals.


\section{The cube-root analog}
\label{sec:cuberoot}

We now consider integrands of the form
\begin{equation}\label{eq:cuberoot-integrand}
  J = \int F(t)\,\frac{\d t}{\sqrt[3]{R(t)}}, \qquad F\in\C(t),\ R\in\C[t].
\end{equation}
The natural geometric setting is the cyclic cubic cover $X_R \colon y^3 = R(t)$.

\subsection{Genus arithmetic}

A direct application of Riemann--Hurwitz to the projection $X_R\to\mathbb{P}^1$, $(t,y)\mapsto t$, of degree $3$ yields the genus of $X_R$ for $R$ of degree $n$ with simple roots:
\[
  g(X_R) \;=\;
  \begin{cases}
    n-1 & \text{if } n\not\equiv 0\pmod 3, \\
    n-2 & \text{if } n\equiv 0\pmod 3.
  \end{cases}
\]
The case relevant to us is genus $1$:
\begin{center}
\begin{tabular}{c|ccccc}
  $\deg R$ & 1 & 2 & 3 & 4 & 5 \\\hline
  $g(X_R)$ & 0 & 1 & 1 & 3 & 4
\end{tabular}
\end{center}
We will focus on the genus-$1$ cases $\deg R\in\{2,3\}$. The case $\deg R=3$ is the cleanest analog of Goursat's quartic-square-root setting (also genus $1$), and we treat it first; the case $\deg R = 2$ may be reduced to $\deg R = 3$ by viewing one root as lying at infinity.

\subsection{Order-3 M\"obius substitutions}

Suppose $R$ is cubic with simple roots $\{a,b,c\}\subset\C$.

\begin{lemma}\label{lem:cyclic-mobius}
There exists a unique M\"obius transformation $S\in\PGL_2(\C)$ of order exactly $3$ sending $a\mapsto b\mapsto c\mapsto a$. Its inverse $S^{-1}=S^2$ sends $a\mapsto c\mapsto b\mapsto a$. The two fixed points $\alpha,\beta\in\mathbb{P}^1(\C)$ of $S$ are not among $\{a,b,c\}$.
\end{lemma}

\begin{proof}
A M\"obius transformation is determined by its action on three points, so $S$ exists and is unique. To see $S$ has order exactly $3$: it is non-trivial since $S(a)=b\neq a$, and $S^3(a)=S^2(b)=S(c)=a$ together with the analogous identities on $b,c$ show that $S^3$ fixes three distinct points $\{a,b,c\}$, hence $S^3=\mathrm{id}$. So $\operatorname{ord}(S)$ divides $3$, and is therefore exactly $3$.

For the fixed points: if $S$ fixed some $r\in\{a,b,c\}$, then $S(r)=r$ would contradict the cyclic permutation $a\mapsto b\mapsto c\mapsto a$ (which has no fixed point in $\{a,b,c\}$).
\end{proof}

In coordinates adapted to $S$, the cyclic action takes a particularly simple form.

\begin{lemma}\label{lem:cyclic-canonical}
Let $\alpha,\beta$ be the fixed points of $S$ in Lemma~\ref{lem:cyclic-mobius}, and set
\[
  z := \frac{t-\alpha}{t-\beta}.
\]
Then $S$ acts in $z$-coordinates as $z\mapsto \omega z$, where $\omega=e^{2\pi i/3}$, and there exist non-zero constants $c,K\in\C$ with
\begin{equation}\label{eq:R-as-z3-K}
  R(t(z))\,(1-z)^3 \;=\; c\,(z^3 - K).
\end{equation}
\end{lemma}

\begin{proof}
The substitution $z=(t-\alpha)/(t-\beta)$ sends $\alpha\mapsto 0$ and $\beta\mapsto\infty$, conjugating $S$ to a M\"obius transformation of $z$ fixing $\{0,\infty\}$, hence multiplication by a scalar $\lambda$. The order condition $\lambda^3=1$ together with $\lambda\neq 1$ forces $\lambda\in\{\omega,\omega^2\}$; relabelling $S\leftrightarrow S^2$ if necessary, we take $\lambda=\omega$.

The inverse change of variable is $t(z)=(\alpha-\beta z)/(1-z)$. As in Step~2 of the proof of Theorem~\ref{thm:goursat1}, multiplying through by $(1-z)^3$ converts $R(t(z))$ into a polynomial of degree $3$ in $z$. Its roots in $z$-coordinates are the images of $a,b,c$ under $r\mapsto(r-\alpha)/(r-\beta)$. Because $S$ acts as multiplication by $\omega$ and permutes $\{a,b,c\}$ cyclically, those images form a single $\omega$-orbit $\{\zeta,\omega\zeta,\omega^2\zeta\}$. Hence
\[
  R(t(z))(1-z)^3 \;=\; c\,(z-\zeta)(z-\omega\zeta)(z-\omega^2\zeta) \;=\; c\,(z^3 - \zeta^3),
\]
and $K=\zeta^3$.
\end{proof}

\subsection{The eigendecomposition}

In $z$-coordinates the cube-radical differential takes a particularly simple form. Writing $t=t(z)=(\alpha-\beta z)/(1-z)$, $\d t=(\alpha-\beta)\,\d z/(1-z)^2$, and using~\eqref{eq:R-as-z3-K}:
\begin{equation}\label{eq:H-defn}
  \frac{F(t)\,\d t}{R(t)^{1/3}} \;=\; \frac{(\alpha-\beta)\,F(t(z))\,\d z\,/(1-z)^2}{c^{1/3}(z^3-K)^{1/3}/(1-z)} \;=\; H(z)\,\frac{\d z}{(z^3-K)^{1/3}},
\end{equation}
where
\begin{equation}\label{eq:H-formula}
  H(z) \;:=\; \frac{(\alpha-\beta)\,F(t(z))}{c^{1/3}\,(1-z)}.
\end{equation}
Here we have chosen a specific cube root $c^{1/3}$ once and for all; different choices alter $H$ and the radical $(z^3-K)^{1/3}$ by mutually cancelling third roots of unity.

The substitution $\sigma\colon z\mapsto \omega z$ generates a cyclic group of order $3$ acting on $\C(z)$. Under this action, $\C(z)$ decomposes as a direct sum of three eigenspaces:
\[
  \C(z) \;=\; V_0 \oplus V_1 \oplus V_2,\qquad V_k = \{f \in \C(z) : f(\omega z) = \omega^k f(z)\}.
\]
The eigenspaces are characterized by
\[
  V_0 = \C(z^3),\quad V_1 = z\cdot\C(z^3),\quad V_2 = z^2\cdot\C(z^3),
\]
and the projection onto $V_k$ is
\begin{equation}\label{eq:projection-Vk}
  P_k f \;=\; \frac{1}{3}\bigl(f(z) + \omega^{-k}f(\omega z) + \omega^{-2k}f(\omega^2 z)\bigr).
\end{equation}

We can therefore write $H = H_0 + H_1 + H_2$ uniquely with $H_k = P_k H \in V_k$, and each piece contributes a separate integral
\[
  J_k := \int H_k(z)\,\frac{\d z}{(z^3-K)^{1/3}}.
\]
The integrals $J_0,J_1,J_2$ behave very differently from one another, and this is the main phenomenon.

\subsection{The main theorem}

\begin{theorem}[Cube-root pseudo-elementarity]\label{thm:cubic-goursat}
Let $R\in\C[t]$ be a cubic polynomial with simple roots, let $S\in\PGL_2(\C)$ be the order-$3$ M\"obius transformation cyclically permuting the roots of $R$, and let $H\in\C(z)$ be the function defined by~\eqref{eq:H-formula} for $F\in\C(t)$. Decompose $H=H_0+H_1+H_2$ as above. Then:
\begin{enumerate}[(i)]
  \item The integrals $J_0$ and $J_2$ are elementary.
  \item The integral $J_1$ equals $\frac{1}{3}\int \varphi_1(x)\,\d x/\sqrt[3]{x(x-K)}$ for $\varphi_1$ the rational function with $H_1(z)=z\,\varphi_1(z^3)$, and is generically not elementary.
\end{enumerate}

In particular, the integral $\int F(t)\,\d t/\sqrt[3]{R(t)}$ is elementary if and only if $H_1\equiv 0$.
\end{theorem}

\begin{proof}
For each $k$, write $H_k(z)=z^k\,\varphi_k(z^3)$ with $\varphi_k\in\C(x)$, where we set $x=z^3$. Substituting into $J_k = \int H_k(z)\,\d z/(z^3-K)^{1/3}$ and using $\d x = 3z^2\,\d z$:
\begin{align*}
  J_0 &= \int \varphi_0(z^3)\,\frac{\d z}{(z^3-K)^{1/3}}
       = \int \varphi_0(x)\,\frac{\d x/(3z^2)}{(x-K)^{1/3}}
       = \frac{1}{3}\int \frac{\varphi_0(x)\,\d x}{x^{2/3}(x-K)^{1/3}}, \\[2pt]
  J_1 &= \int z\,\varphi_1(z^3)\,\frac{\d z}{(z^3-K)^{1/3}}
       = \int \varphi_1(x)\,\frac{\d x/(3z)}{(x-K)^{1/3}}
       = \frac{1}{3}\int \frac{\varphi_1(x)\,\d x}{x^{1/3}(x-K)^{1/3}}, \\[2pt]
  J_2 &= \int z^2\,\varphi_2(z^3)\,\frac{\d z}{(z^3-K)^{1/3}}
       = \int \varphi_2(x)\,\frac{\d x/3}{(x-K)^{1/3}}
       = \frac{1}{3}\int \frac{\varphi_2(x)\,\d x}{(x-K)^{1/3}}.
\end{align*}
The radicand in $J_0$ is $x^2(x-K)$; in $J_1$ it is $x(x-K)$; in $J_2$ it is $(x-K)$. We treat each.

\smallskip\noindent
\textbf{(i$_2$): $J_2$ is elementary.}
The radical $(x-K)^{1/3}$ defines a curve isomorphic to the affine line via $x=K+w^3$, after which $\d x = 3w^2\,\d w$ and
\[
  J_2 = \int \varphi_2(K+w^3)\,w^2\,\d w/w = \int w\,\varphi_2(K+w^3)\,\d w,
\]
which is the integral of a rational function in $w$ and hence elementary.

\smallskip\noindent
\textbf{(i$_0$): $J_0$ is elementary.}
The radicand $x^2(x-K)$ has a \emph{double} root at $x=0$ and a \emph{simple} root at $x=K$. The cube-radical curve $Y_0\colon y^3 = x^2(x-K)$ has genus zero. To see this concretely, set $w = y/x$; then $w^3 = y^3/x^3 = (x-K)/x = 1 - K/x$, hence
\begin{equation}\label{eq:Y0-param}
  x \;=\; \frac{K}{1-w^3},\qquad y \;=\; w x \;=\; \frac{Kw}{1-w^3}
\end{equation}
parametrizes $Y_0$ rationally. Under this parametrization,
\[
  \d x = \frac{3Kw^2}{(1-w^3)^2}\,\d w, \qquad y = \frac{Kw}{1-w^3}, \qquad \frac{\d x}{y} = \frac{3w}{1-w^3}\,\d w,
\]
so $J_0 = \int \frac{\varphi_0(x)}{3y}\d x$ becomes $\int \frac{\varphi_0(K/(1-w^3))\,w}{1-w^3}\,\d w$, the integral of a rational function in $w$, hence elementary.

\smallskip\noindent
\textbf{(ii): $J_1$ is generically transcendental.}
The radicand $x(x-K)$ has two simple roots, and the cube-radical curve $Y_1\colon y^3 = x(x-K)$ has genus $1$ (Riemann--Hurwitz: degree-$3$ map to $\mathbb{P}^1$ ramified at $0,K,\infty$, all with $e=3$, gives $2g-2 = -6 + 3\cdot 2 = 0$, so $g=1$). The differential
\[
  \omega_1 := \frac{\d x}{\sqrt[3]{x(x-K)}}
\]
is holomorphic everywhere on $Y_1$: at each ramification point $P_0,P_K,P_\infty$, a local uniformizer $\tau$ satisfies $x\sim\text{const}\cdot\tau^3$ near $P_0$ and $P_K$ (and $1/x\sim \text{const}\cdot\tau^3$ near $P_\infty$), and an explicit local computation shows $\omega_1$ extends to a regular section of the canonical bundle. Therefore $\omega_1$ is a non-zero element of the one-dimensional space $H^0(Y_1,\Omega^1)$.

By the algebraic version of Liouville's theorem (Theorem~\ref{thm:liouville-alg}, applied to $K=\C(x,(x(x-K))^{1/3})$, together with the result quoted in Section~\ref{sec:liouville} that holomorphic differentials on positive-genus curves never have elementary primitives), the integral $\int \omega_1$ is not elementary. The same conclusion holds for $\int \varphi_1(x)\omega_1$ for generic $\varphi_1\in\C(x)$, since the additional rational factor introduces only further poles whose residues at most produce additional logarithm terms, none of which can absorb the global topological obstruction. (A detailed accounting, in the spirit of the residue calculus on $Y_1$, identifies a precise sufficient condition on $\varphi_1$ for $J_1$ to remain non-elementary; we refer to~\cite{Trager1984,Bronstein2005} for the algorithmic version.)

Combining (i$_0$), (i$_2$), and (ii), the integral $\int F\,\d t/\sqrt[3]{R}$ is elementary if and only if $H_1\equiv 0$.
\end{proof}

\begin{remark}
The dichotomy in Theorem~\ref{thm:cubic-goursat} has \emph{no} analog in the square-root case. Under a $\Z/2$-action there are only two eigenvalues $\pm 1$; the trivial-character component is killed by the orbit-sum (giving Theorem~\ref{thm:goursat2}), and the unique non-trivial component descends elementarily (Theorem~\ref{thm:goursat1}). For the $\Z/3$-action there are three eigenvalues $1,\omega,\omega^2$, and the middle one is genuinely obstructive: it descends only to the cube-root analog of a Weierstrass elliptic curve, $y^3 = x(x-K)$.
\end{remark}

\begin{remark}[Degree-$2$ radicand]\label{rmk:degree-2}
For $R$ of degree $2$, the cube-radical curve $y^3 = R(t)$ still has genus $1$, and Theorem~\ref{thm:cubic-goursat} applies without modification. The proof of Lemma~\ref{lem:cyclic-canonical} extends uniformly: when $\deg R = 2$, only two factors $t(z) - r_i$ contribute and produce a $1/(1-z)^2$ rather than $1/(1-z)^3$, but multiplying through by $(1-z)^3$ still gives a polynomial of degree $3$ in $z$. The "extra" factor $(1-z)$ encodes the implicit ramification point at $\infty$ of the projection $X_R\to\mathbb{P}^1$.

Two structural features specific to $\deg R = 2$ deserve emphasis.

\medskip\noindent
\textbf{(a) $K = 1$ is forced.} Because $S$ cyclically permutes the two finite roots and the point at $\infty$, the three roots in $z$-coordinates form a $\langle\omega\rangle$-orbit on $\mathbb{P}^1(\C)$ containing $z(\infty) = 1$, hence are exactly $\{1,\omega,\omega^2\}$. Therefore
\[
  R(t(z))(1-z)^3 \;=\; c\,(z^3 - 1).
\]
For cubic $R$, by contrast, $K$ depends on the configuration of the roots and may be any non-zero complex number.

\medskip\noindent
\textbf{(b) The fixed points of $S$ are non-real complex conjugates (when $R\in\R[t]$).} A direct computation shows that the unique cyclic M\"obius transformation permuting $\{r_1, r_2, \infty\}$ as $r_1 \to r_2 \to \infty \to r_1$ is
\[
  S(t) \;=\; \frac{r_1\,t - (r_1^2 - r_1 r_2 + r_2^2)}{t - r_2},
\]
with fixed-point equation $t^2 - (r_1+r_2)\,t + (r_1^2 - r_1 r_2 + r_2^2) = 0$ and discriminant $-3(r_1-r_2)^2$. Since $r_1\neq r_2$, the discriminant is strictly negative, hence
\[
  \alpha,\beta \;=\; \frac{r_1+r_2}{2} \;\pm\; \frac{i\sqrt{3}}{2}\,(r_1-r_2),
\]
a pair of non-real complex conjugates. In particular, the change of variable $z = (t-\alpha)/(t-\beta)$ is unavoidably non-trivial — there is no degree-$2$ radicand for which it degenerates to $z = t - \alpha$, in contrast to the cubic case (e.g.\ $R = t^3 - 1$).
\end{remark}

The two features above combine to force a uniform conclusion when $F = 1$:

\begin{corollary}[Non-elementarity for quadratic radicand]\label{cor:degree2-nonelementary}
Let $R(t)\in\C[t]$ be a quadratic polynomial with two distinct roots. Then
\[
  \int \frac{\d t}{\sqrt[3]{R(t)}}
\]
is \emph{not} elementary.
\end{corollary}

\begin{proof}
Specialize Theorem~\ref{thm:cubic-goursat} to $F = 1$. By~\eqref{eq:H-formula},
\[
  H(z) \;=\; \frac{(\alpha-\beta)\,F(t(z))}{c^{1/3}(1-z)} \;=\; \frac{c_0}{1-z},\qquad c_0 := \frac{\alpha-\beta}{c^{1/3}}.
\]
By Remark~\ref{rmk:degree-2}(b), both fixed points $\alpha,\beta$ are finite, so $\alpha-\beta\neq 0$ and hence $c_0\neq 0$. Apply the geometric-series identity
\[
  \frac{1}{1-z} \;=\; \frac{1+z+z^2}{1-z^3}
\]
to read off the eigendecomposition:
\[
  H_0 \;=\; \frac{c_0}{1-z^3}, \qquad H_1 \;=\; \frac{c_0\,z}{1-z^3}, \qquad H_2 \;=\; \frac{c_0\,z^2}{1-z^3}.
\]
Since $c_0\neq 0$, we have $H_1\not\equiv 0$, and Theorem~\ref{thm:cubic-goursat} concludes that the integral is not elementary.
\end{proof}

\begin{remark}
The contrast with the square-root case is striking: every integral $\int\d t/\sqrt{R(t)}$ with $\deg R = 2$ is elementary (yielding $\arcsin$, $\operatorname{arccosh}$, or a logarithm depending on the sign pattern), because the curve $y^2 = R(t)$ has genus $0$. The cube-radical curve $y^3 = R(t)$ has genus $1$, and Corollary~\ref{cor:degree2-nonelementary} shows that for $F=1$ the obstruction in Theorem~\ref{thm:cubic-goursat}(ii) is generic — indeed, unavoidable. The point is that under the (unavoidable) substitution $z = (t-\alpha)/(t-\beta)$, the constant function $F(t) = 1$ does not remain constant: it acquires a pole at $z = 1$ (corresponding to $t = \infty$), and that single pole is enough to populate all three eigencomponents of $H$.
\end{remark}

\begin{remark}[Theorem 2 analog]
For $R$ cubic, the full symmetry group of the unordered triple $\{a,b,c\}$ is the symmetric group $S_3$, comprising the cyclic $\Z/3$ subgroup considered above together with three transpositions (M\"obius involutions swapping pairs). Combining the cyclic and transposition characters yields a six-term sum criterion analogous to the Klein four-group condition~\eqref{eq:V4-sum} of Theorem~\ref{thm:goursat2}; we omit the details.
\end{remark}


\section{Worked examples}
\label{sec:examples}

We now illustrate the framework with a sequence of explicit calculations. Throughout, $\omega = e^{2\pi i/3}$.

\subsection{Square-root: a basic Goursat example}

\begin{example}\label{ex:goursat-basic}
Take $R(t) = (t^2-1)(t^2-4)$ and $F(t)=t$. The roots are $\{\pm 1, \pm 2\}$, and the three M\"obius involutions of Lemma~\ref{lem:involutions-pairs} are computed directly:
\[
  S_1(t) = -t,\qquad S_2(t) = 2/t,\qquad S_3(t) = -2/t.
\]
Verification of the $V_4$-sum:
\[
  F + F\!\circ\! S_1 + F\!\circ\! S_2 + F\!\circ\! S_3 \;=\; t + (-t) + \tfrac{2}{t} + \bigl(-\tfrac{2}{t}\bigr) \;=\; 0.
\]
The character decomposition (with the labelling of Theorem~\ref{thm:goursat2}'s proof) gives
\[
  F^{(1)} = 0,\qquad F^{(2)} = \frac{t^2-2}{2t},\qquad F^{(3)} = \frac{t^2+2}{2t}.
\]
Apply Theorem~\ref{thm:goursat1} to $F^{(2)}$ with the involution $S_1(t)=-t$ (fixed points $\alpha=0$, $\beta=\infty$): since $\beta=\infty$ the substitution is just $u=t$, and~\eqref{eq:R-becomes-even} becomes the trivial $R(u) = (u^2-1)(u^2-4)$, already even. Apply $x = u^2$ to obtain
\[
  \int F^{(2)}\,\frac{\d t}{\sqrt{R}} \;=\; \int \frac{t^2-2}{2t}\,\frac{\d t}{\sqrt{R(t)}}
  \;=\; \tfrac{1}{2}\int \frac{(x-2)\,\d x}{2x\sqrt{(x-1)(x-4)}}.
\]
A similar computation handles $F^{(3)}$. Adding the two contributions and simplifying yields
\[
  \int \frac{t\,\d t}{\sqrt{(t^2-1)(t^2-4)}} \;=\; \tfrac{1}{2}\log\!\Bigl(2t^2 - 5 + 2\sqrt{(t^2-1)(t^2-4)}\Bigr) + C,
\]
which one verifies by direct differentiation. \hfill$\square$
\end{example}

\subsection{Cube root: the canonical setting $R(t)=t^3-1$}

For the next three examples, $R(t)=t^3-1$. The roots are $\{1,\omega,\omega^2\}$, and the order-$3$ M\"obius transformation cyclically permuting them is simply $S(t) = \omega t$. The fixed points of $S$ are $\alpha=0,\beta=\infty$, so the substitution $z=(t-\alpha)/(t-\beta)$ degenerates to $z=t$ and no preliminary change of coordinates is required. The constants in~\eqref{eq:R-as-z3-K} are $c=1$, $K=1$, so the canonical radical is $\sqrt[3]{z^3-1}=\sqrt[3]{t^3-1}$.

In this setting, the eigendecomposition of $H(z)=F(z)$ under $z\mapsto\omega z$ is exactly the eigendecomposition of $F$ under multiplication by $\omega$.

\begin{example}[$F=1$: elementary]\label{ex:F1}
$H(z)=1\in V_0$, so $H_1\equiv 0$ and Theorem~\ref{thm:cubic-goursat} predicts elementarity. With $\varphi_0(x)=1$ and $K=1$:
\[
  \int \frac{\d t}{\sqrt[3]{t^3-1}} \;=\; J_0 \;=\; \frac{1}{3}\int \frac{\d x}{\sqrt[3]{x^2(x-1)}}.
\]
Apply the parametrization~\eqref{eq:Y0-param} of $Y_0\colon y^3=x^2(x-1)$:
\[
  x = \frac{1}{1-w^3},\qquad y = \frac{w}{1-w^3},\qquad \frac{\d x}{y} = \frac{3w}{1-w^3}\,\d w.
\]
Hence
\[
  J_0 \;=\; \frac{1}{3}\int \frac{\d x}{y} \;=\; \int \frac{w\,\d w}{1-w^3}.
\]
The remaining integrand has the partial-fraction decomposition
\[
  \frac{w}{1-w^3} \;=\; \frac{w}{(1-w)(1+w+w^2)} \;=\; \frac{A}{1-w} + \frac{Bw+C}{1+w+w^2}
\]
with $A=\tfrac{1}{3}$, $B=\tfrac{1}{3}$, $C=-\tfrac{1}{3}$, giving
\[
  \int \frac{w\,\d w}{1-w^3} \;=\; -\frac{1}{3}\log(1-w) + \frac{1}{6}\log(1+w+w^2) - \frac{1}{\sqrt{3}}\,\arctan\!\Bigl(\frac{2w+1}{\sqrt{3}}\Bigr) + C.
\]
Inverting the parametrization to recover $w$ in terms of $t$: from $x=t^3=1/(1-w^3)$ we get $1-w^3=1/t^3$, hence $w^3 = (t^3-1)/t^3$ and
\[
  w \;=\; \frac{\sqrt[3]{t^3-1}}{t}.
\]
Substituting back yields the explicit antiderivative:
\[
\boxed{
\begin{aligned}
  \int \frac{\d t}{\sqrt[3]{t^3-1}}
   = &-\tfrac{1}{3}\log\!\Bigl(1-\tfrac{\sqrt[3]{t^3-1}}{t}\Bigr)
     + \tfrac{1}{6}\log\!\Bigl(1+\tfrac{\sqrt[3]{t^3-1}}{t} + \tfrac{\sqrt[3]{(t^3-1)^2}}{t^2}\Bigr) \\
   & - \tfrac{1}{\sqrt{3}}\arctan\!\Bigl(\tfrac{1}{\sqrt{3}}\bigl(2\tfrac{\sqrt[3]{t^3-1}}{t} + 1\bigr)\Bigr) + C.
\end{aligned}
}
\]
One verifies this directly: with $w(t) = \sqrt[3]{t^3-1}/t$, a short calculation gives $\d w/\d t = 1/(t^2(t^3-1)^{2/3})$, while $w/(1-w^3) = w\cdot t^3 = t^2 \sqrt[3]{t^3-1}$, so that the chain rule yields $\d/\d t \int w\,\d w/(1-w^3) = w/(1-w^3)\cdot \d w/\d t = 1/\sqrt[3]{t^3-1}$, as required. \hfill$\square$
\end{example}

\begin{example}[$F=t^2$: trivially elementary]\label{ex:Ft2}
$H(z)=z^2\in V_2$, so $H_0=H_1=0$ and only $J_2$ contributes. With $\varphi_2(x)=1$:
\[
  J_2 \;=\; \frac{1}{3}\int \frac{\d x}{(x-1)^{1/3}}\bigg|_{x=t^3} \;=\; \frac{1}{3}\cdot\frac{3}{2}(x-1)^{2/3}\bigg|_{x=t^3} + C \;=\; \frac{1}{2}(t^3-1)^{2/3} + C.
\]
This recovers the obvious identity $\int t^2\,\d t/\sqrt[3]{t^3-1} = \tfrac{1}{2}\bigl(\sqrt[3]{t^3-1}\bigr)^2 + C$, which one would obtain by inspection (the integrand is, up to a factor of $3$, the derivative of $(t^3-1)^{2/3}$). The eigenvalue framework is consistent with --- and gives a different proof of --- this elementary observation. \hfill$\square$
\end{example}

\begin{example}[$F=t$: \emph{not} elementary]\label{ex:Ft}
$H(z)=z\in V_1$, so $H_0=H_2=0$ and only $J_1$ contributes. With $\varphi_1(x)=1$ and $K=1$:
\[
  \int \frac{t\,\d t}{\sqrt[3]{t^3-1}} \;=\; J_1 \;=\; \frac{1}{3}\int \frac{\d x}{\sqrt[3]{x(x-1)}}.
\]
The right-hand integral is the canonical period integral of the genus-$1$ curve $y^3=x(x-1)$, which is a non-elementary Abelian integral by the discussion in Theorem~\ref{thm:cubic-goursat}(ii).

Concretely, this is a cube-root analog of an incomplete elliptic integral. It is connected to the Chowla--Selberg-type periods of $y^3=x(x-1)$ and can be expressed in terms of hypergeometric functions, but no expression in elementary functions exists. \hfill$\square$
\end{example}

The contrast among the three examples~\ref{ex:F1}--\ref{ex:Ft} is striking: $\int \d t/\sqrt[3]{t^3-1}$ is elementary (yet non-trivially so), $\int t^2\,\d t/\sqrt[3]{t^3-1}$ is elementary (and trivially), and $\int t\,\d t/\sqrt[3]{t^3-1}$ is genuinely transcendental --- a definite ``elliptic-like'' integral living on $y^3=x(x-1)$.

\subsection{A combined integrand}

\begin{example}\label{ex:combined}
Take $R(t)=t^3-1$ and $F(t) = 1 + 5t^2 - 7\,t/(t^3+1)$. Decompose $F$:
\[
  F = 1 + 5t^2 + \widetilde{F}(t), \qquad \widetilde{F}(t) := -\tfrac{7t}{t^3+1}.
\]
The first two terms are in $V_0$ and $V_2$ respectively. The third term $\widetilde{F}(t)$ is of the form $t\cdot\varphi(t^3)$ with $\varphi(x)=-7/(x+1)$, hence in $V_1$. By Theorem~\ref{thm:cubic-goursat}, the integral
\[
  \int F\,\frac{\d t}{\sqrt[3]{t^3-1}}
\]
is elementary if and only if $\widetilde{F} \equiv 0$, which it is not. Hence the integral splits as
\[
  \underbrace{\int \frac{\d t}{\sqrt[3]{t^3-1}}}_{\text{Example~\ref{ex:F1}}} + \underbrace{5\int \frac{t^2\,\d t}{\sqrt[3]{t^3-1}}}_{\text{Example~\ref{ex:Ft2}}} - \tfrac{7}{3}\int \frac{\d x}{(x+1)\sqrt[3]{x(x-1)}},
\]
where the last integral is a non-elementary Abelian integral on $Y_1\colon y^3=x(x-1)$. The framework therefore not only detects elementarity but also \emph{quantifies the obstruction} as a specific Abelian integral on a known curve.
\hfill$\square$
\end{example}

\subsection{A non-canonical $S$: shifted cube}

\begin{example}\label{ex:non-canonical-S}
Consider $R(t)=(t-1)(t-2)(t-3)$ and the order-$3$ M\"obius transformation $S$ with $S(1)=2$, $S(2)=3$, $S(3)=1$. Solving the three linear conditions gives
\[
  S(t) = \frac{5t-13}{3t-7}, \qquad \alpha,\beta = 2 \pm \frac{i}{\sqrt 3}.
\]
A direct check confirms $S(1)=-8/-4=2$, $S(2)=-3/-1=3$, $S(3)=2/2=1$, and the fixed points are the roots of $3t^2-12t+13=0$.

The substitution $z=(t-\alpha)/(t-\beta)$ brings $S$ into canonical form $z\mapsto\omega z$ and produces, by Lemma~\ref{lem:cyclic-canonical}, a polynomial $c(z^3-K)$ for explicit $c$ and $K$. One can then test any rational $F$ by computing its eigencomponents in $z$-coordinates. For $F$ chosen to lie entirely in $V_0$ or $V_2$ in the new coordinates, one obtains explicit elementary antiderivatives via the procedure of Theorem~\ref{thm:cubic-goursat}; for $F$ with non-vanishing $V_1$-component, the integral involves a non-elementary period of the genus-$1$ curve $y^3 = x(x-K)$.
\hfill$\square$
\end{example}

\subsection{A degenerate cubic: degree-$2$ radicand}\label{ex:cubic-degree-2}

\begin{example}
We illustrate Corollary~\ref{cor:degree2-nonelementary} concretely with $R(t) = t^2 - 1$, viewed as a cubic with one root at $\infty$ (cf.\ Remark~\ref{rmk:degree-2}). Here $r_1 = 1$, $r_2 = -1$, so the formula of Remark~\ref{rmk:degree-2}(b) gives
\[
  S(t) \;=\; \frac{1\cdot t - (1 + 1 + 1)}{t - (-1)} \;=\; \frac{t-3}{t+1},
\]
with fixed-point equation $t^2 + 3 = 0$ and so $\alpha = i\sqrt{3}$, $\beta = -i\sqrt{3}$. A direct check confirms $S(1) = -1$, $S(-1) = \infty$, $S(\infty) = 1$.

Substituting $z = (t-\alpha)/(t-\beta)$ with inverse $t(z) = i\sqrt{3}\,(1+z)/(1-z)$:
\[
  R(t(z))\,(1-z)^3 \;=\; \bigl[t(z)^2 - 1\bigr](1-z)^3 \;=\; 4(z^3 - 1),
\]
confirming $K = 1$ as predicted by Remark~\ref{rmk:degree-2}(a), with $c = 4$. The function~\eqref{eq:H-formula} for $F = 1$ is
\[
  H(z) \;=\; \frac{2i\sqrt{3}}{\sqrt[3]{4}\,(1-z)} \;=\; \frac{c_0}{1-z},\qquad c_0 := \frac{2i\sqrt{3}}{\sqrt[3]{4}}\neq 0.
\]
By Corollary~\ref{cor:degree2-nonelementary},
\[
  \boxed{\;\int \frac{\d t}{\sqrt[3]{t^2-1}} \quad\text{is \emph{not} elementary.}\;}
\]

By Theorem~\ref{thm:cubic-goursat}(ii), the obstruction takes the explicit form
\[
  J_1 \;=\; \frac{1}{3}\int \frac{\varphi_1(x)\,\d x}{\sqrt[3]{x(x-1)}},\qquad \varphi_1(x) \;=\; \frac{c_0}{1-x},
\]
a non-elementary Abelian integral on the cube-root elliptic curve $y^3 = x(x-1)$. The integrals $J_0$ and $J_2$ contribute elementary pieces (with $\arctan$ and $\log$ terms via the parametrization~\eqref{eq:Y0-param} and the trivial reduction for $J_2$), but they cannot cancel the transcendental $J_1$.

This is consistent with the fact that symbolic-integration systems return the hypergeometric expression $\int\d t/\sqrt[3]{t^2-1} = t\,\mathrm{e}^{-i\pi/3}\,{}_2F_1\bigl(\tfrac{1}{3},\tfrac{1}{2};\tfrac{3}{2};\, t^2\bigr) + C$, which has no closed form in elementary functions.

The contrast with the square-root case $\int\d t/\sqrt{t^2-1} = \operatorname{arccosh}(t) + C$ is geometric: the curve $y^2 = t^2-1$ has genus $0$ (it is a smooth conic), but $y^3 = t^2-1$ has genus $1$. \hfill$\square$
\end{example}

\section{Concluding remarks}
\label{sec:remarks}

\subsection*{Higher radicals.}
Theorem~\ref{thm:cubic-goursat} extends to $n$-th roots, considering integrands $\int F(t)\,\d t/\sqrt[n]{R(t)}$ together with order-$n$ M\"obius transformations $S$ cyclically permuting roots of $R$. After the substitutions $z = (t-\alpha)/(t-\beta)$ and $x = z^n$, the differential decomposes into $n$ eigencomponents under $z\mapsto\zeta z$ (with $\zeta = e^{2\pi i/n}$), the $k$-th of which lands on the curve
\[
  Y_k\colon\; y^n \;=\; x^{n-1-k}(x-K), \qquad k = 0, 1, \dots, n-1.
\]
A direct application of Riemann--Hurwitz to the projection $Y_k\to\mathbb{P}^1$, $(x,y)\mapsto x$, gives
\begin{equation}\label{eq:gen-n-genus}
  g(Y_k) \;=\; \frac{n+1\;-\;\gcd(n,k)\;-\;\gcd(n,k+1)}{2}.
\end{equation}
Since $\gcd(n,k)$ and $\gcd(n,k+1)$ are coprime divisors of $n$ (as $k$ and $k+1$ are coprime), their sum equals $n+1$ if and only if one of them equals $n$, which occurs precisely for $k = 0$ or $k = n-1$. Therefore
\[
  g(Y_k) = 0 \;\iff\; k\in\{0, n-1\},
\]
and every interior $k\in\{1,\dots,n-2\}$ gives a curve of strictly positive genus. The two outermost eigencomponents always descend to elementary integrals --- via the rational parametrizations $x = K/(1-w^n)$ and $x = K + w^n$ for $Y_0$ and $Y_{n-1}$ respectively --- while the interior components are generically transcendental.

\medskip
For small $n$:

\smallskip\noindent
$\bullet$ $\boxed{n=2}$ \emph{(no interior $k$).} Both eigencomponents descend to genus-$0$ curves; recovering Goursat's theorem (Theorem~\ref{thm:goursat1}).

\smallskip\noindent
$\bullet$ $\boxed{n=3}$ \emph{(one interior $k=1$).} The unique obstruction is $Y_1\colon y^3 = x(x-K)$ of genus $1$, the cube-root analog of a Weierstrass elliptic curve --- this is Theorem~\ref{thm:cubic-goursat}.

\smallskip\noindent
$\bullet$ $\boxed{n=4}$ \emph{(two interior $k\in\{1,2\}$, both genus $1$).} A particularly clean rigidity emerges. For $k=1$, the substitution $u = y^2/x$ on $Y_1\colon y^4 = x^2(x-K)$ gives $u^2 = x-K$, hence $x = u^2 + K$ and
\[
  y^2 \;=\; u\cdot x \;=\; u(u^2 + K) \;=\; u^3 + Ku,
\]
the \emph{lemniscate elliptic curve} with $j$-invariant $1728$. For $k=2$, $Y_2\colon y^4 = x(x-K)$ carries the order-$4$ automorphism $y\mapsto iy$, which descends to an order-$4$ automorphism of its elliptic quotient and therefore forces $j = 1728$. The two obstructions thus both reduce, over $\C$, to the unique elliptic curve with complex multiplication by $\Z[i]$ --- the rigidity is forced by the underlying $\Z/4$-symmetry of the cover.

\smallskip\noindent
$\bullet$ $\boxed{n=5}$ \emph{(three interior $k\in\{1,2,3\}$, all genus $2$).} By~\eqref{eq:gen-n-genus} all three interior curves $Y_1, Y_2, Y_3$ have genus $2$, giving three independent transcendental obstructions on hyperelliptic curves.

\smallskip\noindent
$\bullet$ $\boxed{n\geq 6}$ The genera of the $Y_k$ vary non-monotonically with $k$ and the obstruction curves are no longer pairwise isomorphic. For example at $n = 6$ the interior values $k\in\{1,4\}$ give genus $2$ and $k\in\{2,3\}$ give genus $1$ (cube-root and square-root elliptic respectively), a mixed picture lacking the rigidity present at $n\in\{3,4\}$.

\subsection*{Theorem 2 analog for $n=3$.}
For $R$ cubic, the full automorphism group of the unordered triple of roots is $S_3$, with character table comprising the trivial, sign, and $2$-dimensional irreducible representations. A complete cube-root analog of Theorem~\ref{thm:goursat2} should sum over all of $S_3$ (including transpositions), and the obstruction --- which corresponds in the cyclic case to a single eigenspace --- now correponds to the $2$-dimensional irreducible component. Working this out and matching to the cube-root analog of the Klein four-group decomposition is straightforward but notationally heavier.

\subsection*{Algorithmic perspective.}
The Risch--Trager algorithm~\cite{Risch1969,Trager1984,Bronstein2005} provides a uniform decision procedure for elementary integration over algebraic function fields, including cube-radical integrands. In principle it subsumes both Goursat's theorem and our cube-root analog. In practice, however, the Risch--Trager output is usually opaque: it expresses an antiderivative in a specific algebraic form without revealing the underlying geometric structure. The reductions of Theorems~\ref{thm:goursat1} and~\ref{thm:cubic-goursat} are valuable precisely because they preserve and explain the structure, and they often produce simpler closed forms than the algorithmic output.

\subsection*{Differential Galois theory.}
At a higher level, the question of which integrals are elementary is governed by the differential Galois group of the relevant Picard--Vessiot extension. For curves of positive genus, the Galois-theoretic obstructions to elementarity are governed by the Albanese variety and the structure of the divisor class group~\cite{Bertrand1990}. Goursat's theorem and its cube-root analog can be viewed as explicit calculations of these obstructions in the presence of additional curve automorphisms, and as concrete instances of the general principle that the more symmetry a curve has, the more likely a given Abelian integral is to degenerate to elementary functions.



\begin{thebibliography}{99}

\bibitem{Bertrand1990}
D.~Bertrand,
\emph{Extensions de $\mathcal{D}$-modules et groupes de Galois diff\'erentiels},
in: \emph{p-adic analysis (Trento, 1989)}, Lecture Notes in Math., vol.~1454, Springer, 1990, pp.~125--141.

\bibitem{Bronstein2005}
M.~Bronstein,
\emph{Symbolic Integration I: Transcendental Functions},
2nd ed., Algorithms and Computation in Mathematics, vol.~1, Springer, 2005.

\bibitem{Goursat1887}
\'E.~Goursat,
\emph{Note sur quelques int\'egrales pseudo-elliptiques},
Bulletin de la S.~M.~F., \textbf{15} (1887), 106--120.

\bibitem{Risch1969}
R.~H.~Risch,
\emph{The problem of integration in finite terms},
Trans.\ Amer.\ Math.\ Soc.\ \textbf{139} (1969), 167--189.

\bibitem{Ritt1948}
J.~F.~Ritt,
\emph{Integration in Finite Terms},
Columbia University Press, 1948.

\bibitem{Rosenlicht1972}
M.~Rosenlicht,
\emph{Integration in finite terms},
Amer.\ Math.\ Monthly \textbf{79} (1972), 963--972.

\bibitem{Trager1984}
B.~M.~Trager,
\emph{Integration of algebraic functions},
Ph.D.\ thesis, MIT, 1984.

\end{thebibliography}

\end{document}
